% New post-checkpoint-13 consequence. Full root/B/C/D proof audits passed;
% fresh complete Taylor primary passage audits by root and D.
\begin{corollary}[Proved: symplectic and orthogonal simple quotients]
\label{D:symplectic-orthogonal-simple-product}
Let $Q$ and $p$ satisfy the hypotheses of
Proposition~\ref{D:bounded-scalar-projective-growth}.
For $q=p^f$, consider the following nonabelian simple groups, with
$r$ denoting the dimension of the indicated natural matrix representation:
\[
\begin{array}{c|c|l}
 S&r&\text{range}\\ \hline
 \operatorname{PSp}_{2m}(q)&2m&m\ge1,
       \ (m,q)\notin\{(1,2),(1,3),(2,2)\},\\
 \mathrm{P}\Omega_{2m+1}(q)&2m+1&q\text{ odd},\ m\ge1,
       \ (m,q)\ne(1,3),\\
 \mathrm{P}\Omega^{+}_{2m}(q)&2m&m\ge3,\\
 \mathrm{P}\Omega^{-}_{2m}(q)&2m&m\ge2.
\end{array}                                                   \tag{D.SO.1}
\]
There is a constant $C_Q$, independent of $m,f$ and the row, such that
\[
 \#\{\operatorname{Hom}_{\rm cont}(Q,S)
          \text{ modulo ambient }\operatorname{PGL}_r(\mathbf F_q)
          \text{-conjugacy}\}\le q^{C_Qr^2}.                  \tag{D.SO.2}
\]
The quotient of $Q$ defined by all distinct simple quotient kernels
of these types is topologically finitely generated. There is no
bound on their Lie ranks.

In particular this conclusion applies to the original-boundary tame
group of a normal proper surface over an algebraically closed field
of characteristic $p$. It also applies to the fundamental group of
a proper scheme of dimension at most two after removing a finite
set of codimension at least two in every irreducible component it meets.
It also applies to the puncture of a complete excellent normal local
domain of dimension at least two in equicharacteristic $p$, with
algebraically closed residue field.
\end{corollary}

\begin{proof}
We first verify the precise finite-group inputs. The natural cover
\[
 \operatorname{Sp}_{2m}(q)\longrightarrow
                      \operatorname{PSp}_{2m}(q)
\]
is a quotient by its scalar matrices, which are precisely $I$ and
$-I$. Its kernel therefore has order $\gcd(2,q-1)$.
For an orthogonal geometry $(V,Q_V)$, we use
\[
 \Omega(V)=O(V)'\subset\operatorname{GL}(V),\qquad
 \mathrm{P}\Omega(V)
     =\Omega(V)/(\Omega(V)\cap\mathbf F_q^\times I).
\]
A scalar isometry has square one, since
$Q_V(\lambda v)=\lambda^2Q_V(v)$ and $Q_V$ is nonzero.
Consequently this scalar kernel also has order at most two,
and is trivial in characteristic two. These are the natural
$r$-dimensional matrix lifts. We use the definitions and scalar
calculations in \citet[pp.~69 and 136--137]{Taylor1992ClassicalB}.
Their kernels have prime-to-$p$ order, including when $p=2$.

Theorem~8.8, p.~74, of the same source gives exactly the three
symplectic exceptions in (D.SO.1). For quadratic forms with
nondegenerate polar form, Theorem~11.48, p.~162, gives simplicity
whenever the dimension is at least three and the Witt index is
positive, except for $\mathrm{P}\Omega_3(3)$ and the
four-dimensional plus type. All the orthogonal ranges displayed
in (D.SO.1) therefore consist of nonabelian simple groups.
The omitted low-dimensional groups cause no unmentioned family:
$\operatorname{PSp}_2(q)=\operatorname{PSL}_2(q)$, and
$\mathrm{P}\Omega^{-}_4(q)=\operatorname{PSL}_2(q^2)$.
In even characteristic the odd-dimensional orthogonal geometry
has a degenerate polar form; Theorem~11.9, pp.~143--144, identifies
$O_{2m+1}(q)$ with $\operatorname{Sp}_{2m}(q)$.
Whenever the latter group is perfect its derived group is already
the symplectic group in the first row. The only exceptional even
characteristic parameters are $(m,q)=(1,2),(2,2)$, so they introduce
no unbounded family. We do not apply Theorem~11.48 to those
degenerate polar forms.

For completeness we prove a uniform comparison with group order.
The exact formulas from pp.~70 and 165--166 of
\citet{Taylor1992ClassicalB} are
\[
\begin{split}
 |\operatorname{PSp}_{2m}(q)|
  &=\frac{q^{2m^2+m}}{\gcd(2,q-1)}
                   \prod_{i=1}^{m}(1-q^{-2i}),\\
 |\mathrm{P}\Omega_{2m+1}(q)|
  &=\frac{q^{2m^2+m}}2
                   \prod_{i=1}^{m}(1-q^{-2i})\quad(q\text{ odd}),\\
 |\mathrm{P}\Omega^{\epsilon}_{2m}(q)|
  &=\frac{q^{2m^2-m}}{\gcd(4,q^m-\epsilon)}
       (1-\epsilon q^{-m})
                   \prod_{i=1}^{m-1}(1-q^{-2i}),
       \qquad\epsilon\in\{1,-1\}.
\end{split}                                                   \tag{D.SO.3}
\]
For the orthogonal formulas we also use the spinor-index
calculation and its proof, Theorem~11.51, pp.~164--165.
The exceptional four-dimensional plus group over $\mathbf F_2$
is outside (D.SO.1).

For every $q\ge2$, the elementary inequality
$\prod(1-a_i)\ge1-\sum a_i$, for $0\le a_i\le1$, gives
\[
             \prod_{i=1}^{m}(1-q^{-2i})
                 \ge1-\sum_{i\ge1}2^{-2i}=\frac23.
\]
Thus the first order in (D.SO.3) is at least
$q^{2m^2+m}/3\ge q^{2m^2+m-2}\ge q^{(2m)^2/4}$.
For the second order, $q\ge3$, so it is at least
$q^{2m^2+m-1}\ge q^{(2m+1)^2/8}$ for every $m\ge1$.
For either even-dimensional sign and $m\ge2$,
\[
 \frac{(1-\epsilon q^{-m})}{\gcd(4,q^m-\epsilon)}
       \prod_{i=1}^{m-1}(1-q^{-2i})
       \ge\frac{(3/4)(2/3)}4=\frac18.
\]
Its order is therefore at least
$q^{2m^2-m-3}\ge q^{(2m)^2/8}$, since
$\tfrac32m^2-m-3\ge0$ for $m\ge2$.
In all four rows we have proved
\[
                         |S|\ge q^{r^2/8}.                  \tag{D.SO.4}
\]

Proposition~\ref{D:bounded-scalar-projective-growth}, with
$B=2$ and $\epsilon=1/8$, now proves both assertions for $Q$.
In particular the number of distinct quotient kernels of any
one simple isomorphism type is at most $|S|^{8C_Q}$.
The proposition's uniform generating-tuple argument then applies
to the full product over all these isomorphism types.

For the first geometric application, the second-cohomology
hypothesis follows from
Theorem~\ref{D:surface-tame-cross-characteristic-H2}; the
hereditary quadratic representation bounds follow by finite
normal proper covers from
Theorem~\ref{D:tame-global-quadratic-representation-bound}.
For a possibly nonnormal proper scheme of dimension at most two,
the same normalization descent used in
Corollary~\ref{D:nonnormal-bounded-centre-projective-growth}
applies to the ambient $\operatorname{PGL}_r(\mathbf F_q)$-torsors.
On each fixed connected component of the double fibre product,
a descent isomorphism is determined by one geometric value,
so the total additional number is at most
$|\operatorname{PGL}_r(\mathbf F_q)|^c\le q^{cr^2}$ for a fixed
$c$. Effective finite proper descent
\citep[IX, Theorem~4.7]{C_SGA1} completes this application.
It does not require a second-cohomology theorem for the
nonnormal total scheme.
For the last local application, both abstract hypotheses for its
fundamental group, including the hereditary quadratic bound, have
been verified in
Corollary~\ref{D:local-bounded-scalar-simple-product}.
\end{proof}

\begin{remark}[Scope of the family conclusion]
The scalar kernels of the natural symplectic and orthogonal
lifts are uniformly bounded, independently of their ranks.
The analogous natural unitary lift is treated in
Corollary~\ref{D:bounded-centre-unitary-simple-product}, with
its explicit bound on $\gcd(r,q+1)$. Within these
defining-characteristic classical families, the remaining unbounded
natural centres in linear and unitary groups are not covered
by this argument. No representation-count bound in a coefficient
characteristic different from $p$ is asserted. In particular this
does not settle cross-characteristic Lie families whose orders
are divisible by $p$, or unbounded alternating families.
These conclusions concern products of simple
quotients; they do not supply the required bounds for general
monolithic nonabelian crowns.
\end{remark}
